\chapter{Certification Exam Mapping: DP-203 and DP-700}
\index{DP-203}\index{DP-700}\index{exam mapping}\index{certification}%
This appendix maps the Microsoft data-engineering certification objectives onto the chapters of
this book, then turns the map into a study strategy and a self-assessment checklist. Chapter~23
explains the context: DP-203 (Data Engineering on Microsoft Azure) was retired in 2025 and
\textbf{DP-700} (Implementing Data Engineering Solutions Using Microsoft Fabric) is its successor
credential \cite{microsoft2025dp203,microsoft2025dp700}. The preparation method is identical for
both, and this book's twenty-four weeks were sequenced to cover both exams' domains.

\begin{notebox}
Exam skills-measured documents change. Before booking, download the current version for your
target exam and reconcile it against the tables below---the mapping is durable, the exact
percentages are not.
\end{notebox}

\section{How the Exams Differ}
DP-203 tested the classic Azure data stack directly: ADLS Gen2, Synapse, Databricks, ADF,
Stream Analytics, and the security/monitoring wraparound. DP-700 tests the same engineering
competencies---storage design, ingestion, transformation, orchestration, real-time, security,
optimization---expressed through Microsoft Fabric surfaces: lakehouses, warehouses, eventstreams,
pipelines, notebooks, and Power BI with DirectLake. The cross-cloud boxes throughout this book
(and the consolidated tables in the Cross-Cloud Service Equivalence appendix) are the bridge:
each maps an Azure service to its Fabric counterpart, which is precisely the translation DP-700
demands.

\section{DP-203 Domain Mapping}
\index{DP-203!domains}%
The historical DP-203 domains and their book coverage:

\begin{center}
\footnotesize
\begin{tabular}{@{}p{4.4cm}p{1.6cm}p{6.8cm}@{}}
\toprule
\textbf{DP-203 domain} & \textbf{Weight} & \textbf{Book chapters} \\
\midrule
Design and implement data storage & $\approx$40--45\% & Ch. 1 (ADLS Gen2, tiers, lifecycle), Ch. 2 (Azure SQL/PostgreSQL), Ch. 3 (Cosmos DB), Ch. 5--7 (Synapse dedicated and serverless, lakehouse), Ch. 12 (Delta Lake medallion) \\
Design and develop data processing & $\approx$25--30\% & Ch. 9 (Databricks/Spark), Ch. 10 (ADF, Mapping Data Flows), Ch. 11 (CDC, DMS), Ch. 13--15 (Event Hubs, Stream Analytics, Structured Streaming) \\
Design and implement data security & $\approx$10--15\% & Ch. 21 (zero trust, private endpoints, Key Vault), Ch. 23 (masking, Always Encrypted, Defender, Policy), Ch. 8 (Purview) \\
Monitor and optimize data storage and processing & $\approx$10--15\% & Ch. 6 (distribution, workload management), Ch. 8 (governance telemetry), Ch. 22 (observability, FinOps) \\
\bottomrule
\end{tabular}
\end{center}

If you maintain legacy DP-203-era knowledge for a team or an interview loop, the chapters above
remain the correct study path; only the credential is retired, not the platform.

\section{DP-700 Domain Mapping}
\index{DP-700!domains}%
DP-700 organizes the same competencies around the Fabric platform:

\begin{center}
\footnotesize
\begin{tabular}{@{}p{4.4cm}p{3.6cm}p{4.8cm}@{}}
\toprule
\textbf{DP-700 domain} & \textbf{Book chapters} & \textbf{Fabric surface to learn} \\
\midrule
Implement and manage an analytics solution & Ch. 0--1, 21--24 & Workspace and capacity management, OneLake security, Git integration, deployment pipelines \\
Ingest and transform data & Ch. 9--15 & Pipelines and dataflows gen2, notebooks and Spark jobs, eventstreams, shortcuts, medallion on OneLake \\
Monitor and optimize an analytics solution & Ch. 5--8, 16, 22 & Capacity Metrics app, warehouse and lakehouse optimization, semantic model and DirectLake performance \\
\bottomrule
\end{tabular}
\end{center}

\noindent Fabric-specific surfaces that this Azure book teaches conceptually but that you should
rehearse hands-on in a Fabric trial capacity: OneLake shortcuts (the cross-cloud boxes call them
out under AWS/Azure/GCP storage), DirectLake storage mode (Chapter 16's cross-cloud box), and
eventstreams with Activator (Chapter 14's cross-cloud box).

\section{Objective-to-Chapter Quick Reference}
\index{exam objectives!chapter map}%
Use this table in the final two weeks: find the objective, rebuild the lab from memory, and check
the sketch in the Lab Solution Sketches appendix only after your attempt.

\begin{center}
\footnotesize
\begin{tabular}{@{}p{6.2cm}p{2.2cm}p{4.4cm}@{}}
\toprule
\textbf{Representative objective} & \textbf{Chapters} & \textbf{Evidence you know it} \\
\midrule
Design data partitioning, tiers, redundancy, and lifecycle & 1, 12 & You chose tiers and formats on evidence, not habit \\
Design and implement relational and NoSQL stores & 2, 3 & You can defend partition keys and consistency levels \\
Design star schemas, distributions, and SCD patterns & 5, 6 & You benchmarked distribution choices with numbers \\
Implement serverless lake querying and CETAS & 7 & You cut bytes scanned with Parquet and pruning \\
Build batch pipelines and Mapping Data Flows & 9, 10 & You built idempotent, parameterized pipelines \\
Implement CDC and hybrid data movement & 11 & You gated a CDC stream with a schema contract \\
Implement streaming ingestion and processing & 13--15 & You reason about watermarks and exactly-once \\
Serve data with BI, RLS, and semantic models & 16 & You built a DirectQuery dashboard with dynamic RLS \\
Operationalize ML: features, pipelines, serving, drift & 17--20 & You shipped and monitored a model end-to-end \\
Design security: identities, private networking, secrets & 21, 23 & You built the private-only perimeter and the erasure drill \\
Implement governance: catalog, lineage, quality gates & 8, 22 & You blocked a bad batch with a gate, not a memo \\
Automate with IaC, dbt, and CI/CD & 24 & The pipeline is the only writer to prod \\
\bottomrule
\end{tabular}
\end{center}

\section{Study Strategy}
\index{study strategy}%
Follow the Chapter 23 method:

\begin{enumerate}
    \item \textbf{Weight by weakness, not comfort.} Score yourself per domain; allocate revision
    time by \emph{domain weight $\times$ your weakness}. The failure mode is re-reading what you
    already know.
    \item \textbf{Rebuild, do not re-read.} For every weak domain, rebuild the corresponding
    Thursday lab from memory in a fresh resource group. Fluency under time pressure comes from
    rebuilding, not highlighting.
    \item \textbf{Mock under time.} Take a full-length practice exam in one sitting. Classify
    every miss as a \emph{knowledge gap} (study it), a \emph{misread} (slow down on ``all
    EXCEPT'' questions), or a \emph{judgment gap}---the ``best'' answer among working options is
    usually the one with the least operational toil and the least privilege.
    \item \textbf{Re-mock the weak domains after one week.} Schedule the real exam when two
    consecutive full mocks clear 80\% with no domain below 70\%. The week after Week 24 is the
    right window: material is fresh and the capstone is your scenario practice.
    \item \textbf{Translate to Fabric for DP-700.} For each Azure service in this book, state its
    Fabric surface from memory using the cross-cloud boxes; that translation skill is the core of
    the successor exam.
\end{enumerate}

\section{Self-Assessment Checklist}
\index{self-assessment}%
Check each item only when you can do it unaided in the lab environment, not merely explain it.

\subsection*{Storage and Databases}
\begin{itemize}
    \item[$\square$] Deploy ADLS Gen2 with HNS, lifecycle policy, soft delete, and immutability (Ch. 1)
    \item[$\square$] Choose redundancy (LRS/ZRS/GRS/GZRS) and an access tier from requirements (Ch. 1)
    \item[$\square$] Deploy Azure SQL Serverless with a geo-replica and execute failover (Ch. 2)
    \item[$\square$] Design a Cosmos DB container: partition key, RU budget, consistency level (Ch. 3)
    \item[$\square$] Implement cache-aside with TTL and invalidation reasoning (Ch. 4)
\end{itemize}

\subsection*{Warehousing and Analytics}
\begin{itemize}
    \item[$\square$] Declare the grain, model a star schema, implement SCD2 with MERGE (Ch. 5)
    \item[$\square$] Benchmark hash/round-robin/replicated distributions and read the DMS plan (Ch. 6)
    \item[$\square$] Query the lake with OPENROWSET and materialize with CETAS (Ch. 7)
    \item[$\square$] Explain columnstore, workload groups, and result-set caching (Ch. 6)
\end{itemize}

\subsection*{Ingestion, Transformation, and Orchestration}
\begin{itemize}
    \item[$\square$] Build a Mapping Data Flow and parameterize it for dev/test/prod (Ch. 10)
    \item[$\square$] Configure CDC from a relational source to the lake with a contract gate (Ch. 11)
    \item[$\square$] Orchestrate event-driven Databricks notebooks with failure notification (Ch. 12)
    \item[$\square$] Explain partitions, shuffles, and Catalyst well enough to debug a slow job (Ch. 9)
\end{itemize}

\subsection*{Streaming}
\begin{itemize}
    \item[$\square$] Produce to Event Hubs with partition keys, event IDs, and Capture (Ch. 13)
    \item[$\square$] Write ASA tumbling/hopping windows with a documented late-arrival policy (Ch. 14)
    \item[$\square$] Implement Structured Streaming with watermarks, checkpoints, and idempotent Delta sinks (Ch. 15)
    \item[$\square$] Contrast Service Bus and Event Hubs settlement semantics (Ch. 13)
\end{itemize}

\subsection*{Serving, ML, and GenAI}
\begin{itemize}
    \item[$\square$] Build a DirectQuery report with dynamic RLS and a marked date table (Ch. 16)
    \item[$\square$] Train in Spark, export ONNX, and score with T-SQL PREDICT (Ch. 17)
    \item[$\square$] Build an AML pipeline with a metric gate and model registration (Ch. 18)
    \item[$\square$] Deploy, monitor drift, and trigger evidence-based retraining (Ch. 19)
    \item[$\square$] Assemble a document-extraction and RAG pipeline with PII controls (Ch. 20)
\end{itemize}

\subsection*{Security, Governance, and DataOps}
\begin{itemize}
    \item[$\square$] Build a private-only data service: private endpoint, DNS, no public access (Ch. 21)
    \item[$\square$] Explain managed identity versus service principal secrets, with rotation (Ch. 21)
    \item[$\square$] Run Purview scans with custom classifications and trace lineage (Ch. 8, 22)
    \item[$\square$] Enforce a pre-load quality gate and quarantine path (Ch. 22)
    \item[$\square$] Implement masking, Always Encrypted, and a crypto-shredding erasure drill (Ch. 23)
    \item[$\square$] Ship Terraform + dbt + CI/CD with OIDC and reviewed plans (Ch. 24)
\end{itemize}

\begin{tipbox}
The interview question bank appendix doubles as exam preparation: every question marked
\emph{intermediate} or \emph{advanced} is phrased the way scenario-based exam items are---a
constraint set and a judgment call, not a definition.
\end{tipbox}
